\documentclass[twocolumn,prl,superscriptaddress,showpacs,floatfix]{revtex4-2}

\usepackage{graphicx}
\usepackage{amsmath,amssymb}
\usepackage{hyperref}
\usepackage{booktabs}
\usepackage{xcolor}

\begin{document}

\title{Quantum simulation of coupled-oscillator synchronisation\\
on a 156-qubit superconducting processor}

\author{Miroslav \v{S}otek}
\affiliation{ANULUM, Marbach SG, Switzerland}
\email{protoscience@anulum.li}

\date{\today}

\begin{abstract}
We present a Kuramoto--XY synchronisation study with heterogeneous natural
frequencies, simulator baselines, and ledgered IBM Heron~r2 evidence. A
Rust/PyO3-assisted pipeline enables entanglement entropy, Krylov complexity,
OTOC scrambling, and Floquet discrete-time-crystal diagnostics for systems
of 2--16~qubits. Legacy ibm\_fez artefacts include
Bell, QKD, VQE, ZNE, Trotter, and UPDE rows that must be cited by committed
artefact path rather than as aggregate validation. Finite-size diagnostics are
consistent with Berezinskii--Kosterlitz--Thouless scaling in the small
systems studied, including a periodic-boundary central-charge estimate
$c = 1.04$ at $n = 8$, while larger sizes show finite-size drift. The
same exact-simulation diagnostics show subharmonic Floquet response and
non-ergodic level-spacing behaviour for heterogeneous frequencies. A
hardware-anchored scaling analysis gives a resource boundary for exact
Hilbert-space simulation; it is not a broad quantum-advantage claim because
Rust Kuramoto ODE baselines remain faster for the measured $n \leq 16$
observable-level regime.
All code, data, and figures are open-source at
\url{https://github.com/anulum/scpn-quantum-control}.
\end{abstract}

\maketitle

%% ====================================================================
\section{Introduction}
%% ====================================================================

The Kuramoto model~\cite{Kuramoto1984} describes $N$ coupled oscillators
with natural frequencies $\omega_i$ and coupling matrix $K_{ij}$:
%
\begin{equation}\label{eq:kuramoto}
  \frac{d\theta_i}{dt} = \omega_i
    + \sum_j K_{ij}\sin(\theta_j - \theta_i).
\end{equation}
%
At critical coupling $K_c$ the system undergoes a synchronisation phase
transition characterised by the order parameter
$R = N^{-1}|\sum_k e^{i\theta_k}|$ jumping from zero to a finite value.

The quantum analogue maps this to the XY spin
Hamiltonian~\cite{Pikovsky2001}:
%
\begin{equation}\label{eq:xy}
  H = -\sum_{i<j} K_{ij}(X_i X_j + Y_i Y_j) - \sum_i \omega_i Z_i,
\end{equation}
%
where the $XX + YY$ flip-flop interaction preserves the in-plane ($S^1$)
dynamics while introducing entanglement and quantum tunnelling between
phase configurations. Superconducting-qubit circuits provide a digital
testbed for this spin-model dynamics: each qubit supplies a controllable
two-level oscillator degree of freedom on the Bloch sphere.

Prior quantum simulation of the XY model uses homogeneous frequencies
($\omega_i = \omega$ for all $i$), preserving translational invariance
and the well-characterised BKT universality
class~\cite{Berezinskii1972,Kosterlitz1973}. Theoretical quantum
Kuramoto models with heterogeneous frequencies
exist~\cite{Pikovsky2001,Acebron2005}, and quantum synchronisation
measures have been proposed~\cite{Ameri2015,Galve2013,Ma2020,Giorgi2013}.
The present work contributes a bounded superconducting-processor workflow
for this heterogeneous-frequency setting, with simulator and hardware
claims separated by artefact provenance.

We study the heterogeneous case using a 16-oscillator coupling matrix
$K_{nm} = K_{\mathrm{base}} \cdot \exp(-\alpha|n-m|)$ with
$K_{\mathrm{base}} = 0.45$, $\alpha = 0.3$, and empirically calibrated
nearest-neighbour anchors. The coupling matrix originates from a
hierarchical oscillator model~\cite{Sotek2025} and provides a
non-trivial benchmark with exponential decay, cross-scale boosts, and
heterogeneous frequencies spanning a factor of~1.5.

%% ====================================================================
\section{Methods}
%% ====================================================================

\subsection{Hamiltonian construction}

The XY Hamiltonian~\eqref{eq:xy} is constructed directly in the
computational basis via bitwise flip-flop operations, bypassing
Qiskit's \texttt{SparsePauliOp}:
%
\begin{align}
  H_{k,\, k \oplus \mathrm{mask}_{ij}} &= -2K_{ij}
    \quad\text{when } b_i(k) \neq b_j(k), \\
  H_{kk} &= -\sum_i \omega_i (1 - 2b_i(k)),
\end{align}
%
where $b_i(k)$ is the $i$-th bit of basis state $|k\rangle$ and
$\mathrm{mask}_{ij}$ flips bits $i$ and $j$. Current benchmark artefacts
report bounded Rust/PyO3 kernel-level acceleration for selected tasks
(Table~\ref{tab:speedup}). Earlier unmatched internal timing numbers are
not reported because they have not been reproduced under the current
matched benchmark harness.

\begin{table}[b]
\caption{\label{tab:speedup}%
Measured Rust vs.\ Python/Qiskit/scipy speedups.}
\begin{ruledtabular}
\begin{tabular}{lrl}
Module & Speedup & Method \\
\colrule
Knm build ($n=4$) & $5.89\times$ & Matched current harness \\
Knm build ($n=16$) & $3.71\times$ & Matched current harness \\
OTOC ($n=4$) & $264\times$ & Eigendecomp + rayon \\
Krylov ($n=3$) & $27\times$ & Complex Lanczos \\
Order param.\ ($n=4$) & $6.2\times$ & Batch Pauli \\
\end{tabular}
\end{ruledtabular}
\end{table}

\subsection{Analysis pipeline}

Simulation diagnostics include half-chain entanglement entropy $S(A)$
and Schmidt gap $\Delta_S = \lambda_1 - \lambda_2$ (via exact
diagonalisation), Krylov complexity and Lanczos coefficients $b_n$
(complex Lanczos algorithm~\cite{delCampo2025}), out-of-time-order
correlators $F(t) = \langle W^\dagger(t) V^\dagger W(t) V\rangle$
(eigendecomposition approach~\cite{Maldacena2016}), and Floquet DTC
analysis (stroboscopic order parameter at subharmonic
frequency)~\cite{Else2016}. Level-spacing statistics use the ratio
$\bar{r} = \langle\min(\delta_n, \delta_{n+1}) /
\max(\delta_n, \delta_{n+1})\rangle$~\cite{Oganesyan2007}.

\subsection{Hardware evidence}

The March~2026 \textbf{ibm\_fez} material is retained as a separate
legacy artefact-backed campaign, not as the promoted dataset for the later
DLA/parity claim. Individual Bell, QKD, VQE, ZNE, Trotter, and UPDE
values must cite their committed result artefact or the hardware ledger. The
promoted raw-count IBM dataset is the April~2026 \textbf{ibm\_kingston}
Phase~1 DLA parity campaign, which includes job IDs, counts, integrity checks,
and a count-to-statistic reproduction harness.

\subsection{Statistical and artefact protocol}

Every numerical claim in this overview is tied to one of three evidence
classes. Exact-simulation rows are deterministic within the specified
Hamiltonian, system size, and solver tolerance. Hardware rows are raw-count
observations with job identifiers, shot counts, and committed reducer scripts.
Scaling rows are resource-boundary estimates and are not promoted to broad
quantum advantage claims. This classification is necessary because the paper
combines heterogeneous evidence: Bell/QKD rows, finite-size exact diagnostics,
legacy IBM rows, and later parity-sector campaigns have different statistical
meaning and should not be collapsed into one aggregate success metric.

\begin{table}[b]
\caption{\label{tab:evidenceclasses}%
Evidence classes used in this overview.}
\footnotesize
\begin{tabular}{@{}p{0.22\columnwidth}p{0.28\columnwidth}p{0.40\columnwidth}@{}}
\toprule
Class & Examples & Claim boundary \\
\midrule
Exact simulation & BKT, Krylov, OTOC, DTC & finite-size diagnostics \\
Raw-count hardware & Bell, QKD, DLA rows & backend/window observations \\
Resource boundary & exact crossover & no broad advantage claim \\
Legacy artefact & early Fez rows & cited only by artefact path \\
\bottomrule
\end{tabular}
\end{table}

\begin{table}[t]
\caption{\label{tab:programmeroadmap}%
Programme roadmap linking the overview evidence to the focused companion
papers. The table is a map of claim ownership, not an aggregate validation
score.}
\footnotesize
\begin{tabular}{@{}p{0.20\columnwidth}p{0.24\columnwidth}p{0.46\columnwidth}@{}}
\toprule
Component & Focused paper & Positive knowledge gained \\
\midrule
Exact diagnostics & This overview & finite-size and resource-boundary map \\
Software methods & Rust/VQE methods & reproducible artefact pipeline and kernel bounds \\
DLA parity & Phase~1 DLA parity & real but multifactorial leakage asymmetry \\
FIM term & FIM Hamiltonian & negative design rule for direct digital feedback \\
Reduced Pauli & Phase~3 tomography & channel-localised mechanism boundary \\
Dynamic feedback & S1 feedback & backend-dependent control-response boundary \\
\bottomrule
\end{tabular}
\end{table}

%% ====================================================================
\section{Simulation results}
%% ====================================================================

\subsection{Entanglement at the synchronisation transition}

The Schmidt gap shows a sharp minimum at $K \approx 3.44$ for $n = 8$
(Fig.~\ref{fig:entanglement}), marking the synchronisation transition.
Entanglement entropy saturates at values consistent with the
Calabrese--Cardy scaling $S \sim (c/3)\ln L$ for a $c = 1$
CFT~\cite{Calabrese2004}.

\begin{figure}[t]
\includegraphics[width=\columnwidth]{figures/publication/fig1_entanglement_vs_coupling.pdf}
\caption{\label{fig:entanglement}%
Half-chain entanglement entropy $S(A)$ and Schmidt gap $\Delta_S$
across coupling strength for $n = 2, 3, 4, 6, 8$ oscillators with
heterogeneous frequencies.}
\end{figure}

\subsection{Krylov complexity}

Mean Lanczos coefficient $\langle b_n \rangle$ grows linearly with $K$,
indicating that operator growth rate scales with coupling strength
(Fig.~\ref{fig:krylov}).

\begin{figure}[t]
\includegraphics[width=\columnwidth]{figures/publication/fig2_krylov_vs_coupling.pdf}
\caption{\label{fig:krylov}%
Peak Krylov complexity and mean Lanczos coefficient vs.\ coupling.}
\end{figure}

\subsection{OTOC scrambling}

Strong coupling scrambles $4\times$ faster: $t^* = 0.28$ ($K=4$) vs.\
$t^* = 1.17$ ($K=1$) at $n=8$ (Fig.~\ref{fig:otoc}).

\begin{figure}[t]
\includegraphics[width=\columnwidth]{figures/publication/fig3_otoc_time_traces.pdf}
\caption{\label{fig:otoc}%
OTOC $F(t)$ at sub-critical ($K=1$) and super-critical ($K=4$)
coupling for $n = 4, 6, 8$.}
\end{figure}

\subsection{Floquet discrete time crystal}

The exact-simulation scan shows subharmonic response across the 15 drive
amplitudes tested with heterogeneous frequencies (Fig.~\ref{fig:dtc}).
Within this finite-size diagnostic, frequency disorder does not remove the
Floquet response. The non-ergodic spectrum (Sec.~\ref{sec:mbl}) is a
candidate stabilisation mechanism rather than a hardware-demonstrated DTC
mechanism.

\begin{figure}[t]
\includegraphics[width=\columnwidth]{figures/publication/fig9_floquet_dtc_n3456.pdf}
\caption{\label{fig:dtc}%
Subharmonic ratio $P(\Omega/2)/P(\Omega)$ and mean $R$ vs.\ drive
amplitude $\delta$ for $n = 3, 4, 6$.}
\end{figure}

\subsection{Finite-size scaling}

The BKT ansatz gives $K_c(\infty) \approx 2.20$ in the current
finite-size fit (Fig.~\ref{fig:fss}). The spectral-gap fits at
$n=4,6,8$ are consistent with BKT-like scaling, and the periodic-boundary
central-charge estimate is $c = 1.04$ at $n = 8$ ($c = 1$ expected for
BKT). The failed or drifting larger-size fits are treated as finite-size
limitations, not as evidence of asymptotic universality.

\begin{figure}[t]
\includegraphics[width=\columnwidth]{figures/publication/fig6_finite_size_scaling.pdf}
\caption{\label{fig:fss}%
Critical coupling $K_c(N)$ from spectral gap minimum. BKT ansatz
(solid) vs.\ power law (dashed).}
\end{figure}

%% ====================================================================
\section{Hardware results}
%% ====================================================================

\subsection{Bell test}

Two independent Bell pairs yield $S_{01} = 2.165 \pm 0.02$ and
$S_{23} = 2.188 \pm 0.02$, both violating the classical limit of~2 at
${>}8\sigma$ significance (Fig.~\ref{fig:hardware}).

\subsection{QKD viability}

Quantum bit error rate in matched bases: ZZ 5.5\%, XX 5.8\% --- below
the BB84 security threshold of 11\%. Mismatched basis (ZX): 93.9\%
error, consistent with basis discrimination in this protocol.

\subsection{ZNE stability}

Mean $\langle Z \rangle$ varies by ${<}2\%$ across fold levels 1--9
for both 4-qubit and 8-qubit systems. Richardson extrapolation provides
${<}2\%$ correction to raw values.

The ZNE row is therefore used as a stability diagnostic, not as a strong
mitigation claim. It shows that this shallow legacy observable is not highly
sensitive to the tested folding levels, but it does not validate all later
DLA, FIM, or feedback observables. Those later papers use their own
readout/ZNE lanes when mitigation is part of the claim.

\subsection{16-qubit dynamics}

13 of 16~qubits show $|\langle Z \rangle| > 0.3$ in the legacy UPDE row,
showing that part of the Kuramoto coupling pattern remains visible, while
the 16-qubit circuit is explicitly noise-limited (Fig.~\ref{fig:hardware16}).

\begin{figure}[t]
\includegraphics[width=\columnwidth]{figures/publication/fig10_ibm_hardware.pdf}
\caption{\label{fig:hardware}%
Legacy artefact-backed \texttt{ibm\_fez} hardware campaign from March 2026,
retained for completeness and workflow context rather than promoted as the
later DLA/parity dataset. (a)~Per-qubit $\langle Z \rangle$ heatmap.
(b)~8-qubit expectations show coupling pattern. (c)~QKD QBER: 5.5\%.
(d)~CHSH $S = 2.165 > 2$.}
\end{figure}

\begin{figure}[t]
\includegraphics[width=\columnwidth]{figures/publication/fig14_complete_analysis.pdf}
\caption{\label{fig:hardware16}%
Legacy \texttt{ibm\_fez} analysis companion to Fig.~\ref{fig:hardware}. These
rows are artefact-backed context for the overview, not the promoted evidence
for the later DLA, FIM, Phase~3, or S1 claims. (a)~ZZ correlation matrix.
(b)~Trotter order comparison. (c)~16-qubit per-qubit $\langle Z \rangle$:
alternating pattern visible. (d)~VQE 8-qubit energy landscape.}
\end{figure}

\subsection{Exact-simulation crossover boundary}

The completed qubit-scaling series provides a direct resource boundary
for exact Hilbert-space simulation. Figure~\ref{fig:advantage} combines
the ibm\_fez hardware jobs at $n=4,6,8,10,12,14$ and the 16-qubit UPDE
snapshot with the committed classical baselines. Exact diagonalisation
grows from 0.1~ms at $n=4$ to 26.8~s at $n=12$ and reaches the memory
wall at $n=14$--16 on the 32~GB workstation used for the baseline. A
log-linear fit to the finite exact timings crosses the QPU budget
envelope at $n \approx 11.6$.

This crossing is not a broad computational-advantage claim. It is a
boundary for exact state-vector/Hilbert-space simulation of the
Kuramoto--XY Hamiltonian. The Rust Kuramoto ODE baseline remains in the
millisecond regime through $n=16$, and the hardware data are
noise-limited beyond the shallowest circuits.

The crossover should be read together with the memory boundary. Exact dense
methods encounter a state-vector/Hamiltonian storage wall before the classical
ODE baseline does, and tensor-network or sparse methods can move that boundary
depending on the observable. The figure is therefore useful as an engineering
map for which classical baseline is being challenged, not as a statement that
the hardware has outperformed all classical descriptions of the oscillator
model.

\begin{figure}[t]
\includegraphics[width=\columnwidth]{figures/publication/fig17_quantum_advantage_crossover.pdf}
\caption{\label{fig:advantage}%
Exact-simulation crossover anchored by completed ibm\_fez hardware runs.
Red points and fit show exact diagonalisation wall time; blue squares and
dashed curve show hardware QPU budget estimates for completed runs; green
triangles show the Rust Kuramoto ODE baseline. The crossover at
$n \approx 11.6$ applies only to exact Hilbert-space simulation.}
\end{figure}

\subsection{Ansatz comparison}

The physics-informed Knm ansatz (CZ gates between coupled pairs)
concentrates 42\% of probability in the top bitstring vs.\ 20\% for
TwoLocal and 21\% for EfficientSU2 in this legacy hardware row, supporting
the narrower claim that topology-reflecting circuit design can improve this
benchmark against the two generic ansatze tested.

\subsection{Non-ergodicity diagnostic}\label{sec:mbl}

Level-spacing ratio $\bar{r}$ distinguishes integrable
($\bar{r} \approx 0.386$, Poisson) from chaotic
($\bar{r} \approx 0.530$, GOE)
spectra~\cite{Oganesyan2007}~(Fig.~\ref{fig:mbl}):
%
\begin{itemize}
  \item $n = 4$: crosses Poisson to near-GOE at $K \approx 2.1$.
  \item $n = 6$: stays mostly below GOE; chaos onset only at $K = 8.0$.
  \item $n = 8$: \emph{never reaches GOE} (max $\bar{r} = 0.43$).
\end{itemize}
%
Eigenstate entanglement sits 30--40\% below the thermal (Page)
expectation but grows with $N$, ruling out deep MBL. The correct
characterisation is a \emph{non-ergodic regime} where heterogeneous
frequencies protect the coupling topology from thermal scrambling.

\begin{figure}[t]
\includegraphics[width=\columnwidth]{figures/publication/fig15_mbl_level_spacing.pdf}
\caption{\label{fig:mbl}%
Level-spacing ratio $\bar{r}$ vs.\ coupling. At $n = 8$ the finite-size
diagnostic remains below the GOE reference across the tested grid.}
\end{figure}

%% ====================================================================
\section{Discussion}
%% ====================================================================

\textbf{BKT-like finite-size diagnostics.}
Two diagnostics support a BKT-like interpretation in the small systems
studied: the periodic-boundary central-charge estimate $c = 1.04$ at
$n = 8$ (expected $c = 1$), and spectral-gap essential-singularity fits
$\Delta \sim \exp(-b/\sqrt{K-K_c})$ with $R^2 > 0.96$ for the supported
sizes. Larger-size drift keeps this as a finite-size diagnostic rather
than an asymptotic-universality conclusion.

\textbf{DTC resilience.}
The non-ergodic spectrum provides a plausible mechanism in exact
simulation: heterogeneous frequencies suppress thermalisation trends and
coincide with a stable subharmonic response without requiring deep MBL.

\textbf{Noise budget.}
The 5.5\% QBER and 94\% state preparation fidelity in the retained
legacy rows indicate usable shallow-circuit performance for the tested
Kuramoto--XY circuits. The 16-qubit row remains explicitly noise-limited
(depth $\leq 50$ CX gates).

\textbf{Limitations.}
No quantum advantage is demonstrated; classical solvers outperform at
$n \leq 16$ for observable-level ODE baselines. The scaling crossover in
Fig.~\ref{fig:advantage} is limited to exact Hilbert-space simulation and
should be read as a resource boundary, not as demonstrated broad
advantage. Trotter error at $dt = 0.1$ is significant (sign flip at finer
resolution). The coupling matrix used here originates from a
hierarchical oscillator model~\cite{Sotek2025}; the Kuramoto-to-XY
mapping itself is standard physics and the analysis applies to any
coupling matrix.
This overview also combines material generated at different stages of the
programme. It should therefore be read as a map of the evidence stack rather
than as a single homogeneous experiment. The Phase~1, Phase~3, FIM, and S1
papers provide the stricter individual claim boundaries for the hardware rows
that are most relevant to publication.
At the time of execution, authenticated live QPU access for the hardware rows
in this study was limited to IBM Quantum hardware. The software and artefact
pipeline has provider-neutral routes for other gate-model, trapped-ion,
neutral-atom, photonic, annealing, analogue, and simulator targets, but those
routes are readiness infrastructure rather than live non-IBM execution
evidence. We plan to rerun the same bounded synchronisation payloads on
additional hardware families when compute allocations and provider access
become available.

\section{Provider access and programme boundaries}

The present stack should be read as an IBM-Heron hardware programme with
provider-neutral software readiness, not as a cross-provider universality
claim. IBM Quantum was the authenticated QPU route available for these
campaigns at execution time. The framework already contains provider and
aggregator surfaces for other hardware families, but a software route is not
evidence of live hardware execution. Individual companion papers therefore
state their IBM boundary briefly, while this overview records the programme
level reason: cross-provider reruns are the next evidence class once compute
allocations and access are available. Until then, the positive contribution is
the raw-count, preregistered, reducer-bound methodology and the exact
mechanism boundaries it establishes inside the available IBM hardware window.

%% ====================================================================
\section{Conclusion}
%% ====================================================================

We have presented a bounded quantum hardware study of
coupled-oscillator synchronisation with heterogeneous natural
frequencies. Key results: Bell inequality violation ($S = 2.165$),
sub-threshold QKD error rate (5.5\%), BKT critical coupling
$K_c \approx 2.2$ in the finite-size fit, subharmonic Floquet response in
exact simulation, and a non-ergodic spectral regime at $n = 8$. The
software pipeline combines Python/Qiskit orchestration with Rust/PyO3
kernel acceleration for selected reproducibility and benchmarking tasks.
All code, data, and 17~figures are open-source (AGPL-3.0) at
\url{https://github.com/anulum/scpn-quantum-control}.

%% ====================================================================
\section*{Data availability}
%% ====================================================================

Simulation data, ledgered hardware artefacts, Python modules, Rust engine
source, and publication figures are available at the GitHub repository under
AGPL-3.0. Zenodo archive:
DOI~\href{https://doi.org/10.5281/zenodo.18821929}{10.5281/zenodo.18821929}.

\begin{thebibliography}{20}

\bibitem{Kuramoto1984}
Y.~Kuramoto,
\emph{Chemical Oscillations, Waves, and Turbulence}
(Springer, Berlin, 1984).

\bibitem{Pikovsky2001}
A.~Pikovsky, M.~Rosenblum, and J.~Kurths,
\emph{Synchronization: A Universal Concept in Nonlinear Sciences}
(Cambridge University Press, 2001).

\bibitem{Acebron2005}
J.~A.~Aceb\'on, L.~L.~Bonilla, C.~J.~P\'erez~Vicente, F.~Ritort, and R.~Spigler,
Rev.\ Mod.\ Phys.\ \textbf{77}, 137 (2005).

\bibitem{Berezinskii1972}
V.~L.~Berezinskii,
Sov.\ Phys.\ JETP \textbf{34}, 610 (1972).

\bibitem{Kosterlitz1973}
J.~M.~Kosterlitz and D.~J.~Thouless,
J.\ Phys.\ C \textbf{6}, 1181 (1973).

\bibitem{Calabrese2004}
P.~Calabrese and J.~Cardy,
J.\ Stat.\ Mech.\ P06002 (2004).

\bibitem{Maldacena2016}
J.~Maldacena, S.~H.~Shenker, and D.~Stanford,
J.\ High Energy Phys.\ \textbf{2016}(8), 106.

\bibitem{delCampo2025}
A.~del Campo \emph{et al.},
arXiv:2510.13947 (2025).

\bibitem{Else2016}
D.~V.~Else, B.~Bauer, and C.~Nayak,
Phys.\ Rev.\ Lett.\ \textbf{117}, 090402 (2016).

\bibitem{Oganesyan2007}
V.~Oganesyan and D.~A.~Huse,
Phys.\ Rev.\ B \textbf{75}, 155111 (2007).

\bibitem{Galve2013}
F.~Galve, G.~L.~Giorgi, and R.~Zambrini,
in \emph{Lectures on General Quantum Correlations and their Applications}
(Springer, 2017), pp.~393--420;
see also Phys.\ Rev.\ A \textbf{81}, 062117 (2010).

\bibitem{Ameri2015}
V.~Ameri, M.~Eghbali-Arani, A.~Mari, A.~Farace, F.~Kheirandish,
V.~Giovannetti, and R.~Fazio,
Phys.\ Rev.\ A \textbf{91}, 012301 (2015).

\bibitem{Ma2020}
J.~Ma, Y.~Chen, and Z.~Wang,
arXiv:2005.09001 (2020).

\bibitem{Giorgi2013}
G.~L.~Giorgi, F.~Galve, G.~Manzano, P.~Colet, and R.~Zambrini,
Phys.\ Rev.\ A \textbf{85}, 052101 (2012).

\bibitem{Wiersema2023}
R.~Wiersema \emph{et al.},
arXiv:2309.05690 (2023).

\bibitem{Sotek2025}
M.~\v{S}otek,
\emph{God of the Math --- The SCPN Master Publications},
Zenodo, DOI:~10.5281/zenodo.17419678 (2025).

\bibitem{IBM2026}
IBM Quantum,
ibm\_fez backend specifications,
\url{https://quantum.cloud.ibm.com} (2026).

\end{thebibliography}

\end{document}
